\documentclass[12pt]{article}
\usepackage[utf8]{inputenc}
\usepackage[spanish]{babel}
\usepackage{geometry}
\usepackage{enumitem}
\usepackage{fancyhdr}

\geometry{a4paper, margin=1in}

\title{Examen de Física Básica: Cinemática y Newton - Forma B}
\author{Nombre: \underline{\hspace{8cm}} \quad Matrícula: \underline{\hspace{3cm}}}
\date{Fecha: \underline{\hspace{4cm}}}

\pagestyle{fancy}
\fancyhf{}
\rhead{Examen Forma B}
\lhead{Física Básica}
\cfoot{\thepage}

\begin{document}

\maketitle

\section*{I. SECCIÓN TEÓRICA (4 Puntos)}

\begin{enumerate}[label=\arabic*.]
    \item \textbf{¿Cómo se denomina al ente constituido por un punto de referencia, un conjunto de ejes coordenados y un reloj?}
    \begin{enumerate}[label=\alph*)]
        \item Partícula
        \item Trayectoria
        \item Posición
        \item Marco de referencia
    \end{enumerate}

    \item \textbf{En el Movimiento Rectilíneo Uniforme (MRU), ¿cuál de las siguientes características se mantiene constante?}
    \begin{enumerate}[label=\alph*)]
        \item La aceleración
        \item La posición
        \item La velocidad (magnitud, dirección y sentido)
        \item El tiempo
    \end{enumerate}

    \item \textbf{En el estudio de la Caída Libre, despreciando la resistencia del aire, la aceleración gravitacional (g) cerca de la superficie terrestre tiene una magnitud aproximada de:}
    \begin{enumerate}[label=\alph*)]
        \item 8.9 $m/s^2$ hacia arriba
        \item 9.8 $m/s^2$ hacia abajo
        \item 0 $m/s^2$
        \item 1.63 $m/s^2$ hacia abajo
    \end{enumerate}

    \item \textbf{Según la Segunda Ley de Newton, la aceleración de un cuerpo es:}
    \begin{enumerate}[label=\alph*)]
        \item Inversamente proporcional a la fuerza neta
        \item Directamente proporcional a la masa
        \item Directamente proporcional a la fuerza neta e inversamente proporcional a la masa
        \item Independiente de la masa y la fuerza
    \end{enumerate}
\end{enumerate}

\section*{II. SECCIÓN PRÁCTICA (6 Puntos)}

\begin{enumerate}[label=\arabic*.]
    \item \textbf{Un corredor recorre 120 metros planos en un tiempo de 6.0 segundos. ¿Cuál fue su velocidad media durante la carrera?}
    \begin{enumerate}[label=\alph*)]
        \item 10 m/s
        \item 20 m/s
        \item 60 m/s
        \item 720 m/s
    \end{enumerate}

    \item \textbf{Determine el desplazamiento de un objeto que se mueve con MRUV si su velocidad inicial es de 10.0 m/s y acelera a 4.0 $m/s^2$ durante 3.0 segundos.}
    \begin{enumerate}[label=\alph*)]
        \item 30 m
        \item 18 m
        \item 48 m
        \item 40 m
    \end{enumerate}

    \item \textbf{Un proyectil es lanzado hacia arriba con una velocidad inicial de 25 m/s. ¿Qué velocidad tendrá al alcanzar su punto de máxima altura?}
    \begin{enumerate}[label=\alph*)]
        \item 25 m/s
        \item 9.8 m/s
        \item 0 m/s
        \item -9.8 m/s
    \end{enumerate}

    \item \textbf{¿Cuál es el peso de un objeto cuya masa es de 5.0 kg en la superficie de la Tierra (g = 9.8 $m/s^2$)?}
    \begin{enumerate}[label=\alph*)]
        \item 5.0 N
        \item 0.51 N
        \item 49.0 N
        \item 9.8 N
    \end{enumerate}
\end{enumerate}

\end{document}
